\documentclass[10pt]{article}

\usepackage[a4paper,margin=0.72in]{geometry}
\usepackage[T1]{fontenc}
\usepackage[utf8]{inputenc}
\usepackage{times}
\usepackage{amsmath,amssymb}
\usepackage{booktabs}
\usepackage{longtable}
\usepackage{enumitem}
\usepackage{xurl}
\usepackage[colorlinks=true,linkcolor=black,citecolor=black,urlcolor=blue]{hyperref}

\setlength{\parindent}{0pt}
\setlength{\parskip}{3pt}
\setlist[itemize]{leftmargin=*,nosep}
\setlist[enumerate]{leftmargin=*,nosep}
\renewcommand{\arraystretch}{1.08}

\title{\textbf{Extracted Reference Notes for the CS240 Project}\\
Communication-Constrained Distributed Influence Maximization}
\author{Prepared from the seven PDFs in \texttt{references/}}
\date{}

\begin{document}
\maketitle

\section{Project Context}

The course project studies \textbf{communication-constrained distributed influence
maximization}.  The target problem is to choose a seed set
$S\subseteq V$, $|S|\le k$, that maximizes expected influence spread
$\sigma(S)$ under a stochastic diffusion model, while comparing centralized
greedy, lazy greedy/CELF, and a simple distributed candidate-sharing scheme.
The most useful material from the references is therefore:
\begin{itemize}
\item the formal influence maximization model and the Independent Cascade (IC)
and Linear Threshold (LT) diffusion processes;
\item the monotone submodular structure of $\sigma(S)$ and the
$(1-1/e)$ greedy approximation guarantee;
\item the computational bottleneck of repeated marginal-gain estimation;
\item lazy evaluation as a way to reduce repeated spread evaluations;
\item distributed submodular maximization as the theoretical template for
``local greedy plus global merge'';
\item modern scalable systems that motivate communication budgets, compression,
parallel seed selection, and reverse influence sampling.
\end{itemize}

\section{One-Page Synthesis}

\begin{longtable}{p{0.21\linewidth}p{0.34\linewidth}p{0.36\linewidth}}
\toprule
Reference & Main idea & Direct use in our project \\
\midrule
\endfirsthead
\toprule
Reference & Main idea & Direct use in our project \\
\midrule
\endhead
Nemhauser et al. 1978 & Classical analysis of maximizing a nondecreasing
submodular set function under a cardinality constraint. Greedy achieves
$1-\left(\frac{K-1}{K}\right)^K$, approaching $1-1/e$. &
Provides the theorem behind our greedy guarantee. We can cite it when explaining
why diminishing returns make greedy a principled approximation algorithm. \\
\midrule
Kempe et al. 2003 & Introduces influence maximization under IC/LT models;
proves NP-hardness; proves $\sigma(S)$ is monotone submodular; applies greedy
for a near $63\%$ guarantee. &
Defines our core problem. Provides IC process, objective $\max_{|S|\le k}\sigma(S)$,
NP-hardness, submodularity, and the comparison against degree/centrality
heuristics. \\
\midrule
Leskovec et al. 2007 & Shows many outbreak-detection objectives are submodular
and develops CELF, a cost-effective lazy forward greedy algorithm. &
Supports our lazy greedy implementation. Its priority-queue logic explains why
old marginal gains can serve as upper bounds after the seed set grows. \\
\midrule
Mirzasoleiman et al. 2013 & Proposes GreeDI, a two-stage distributed protocol:
partition data, run local greedy, merge candidates, run greedy again. &
Main template for our distributed method. We can adapt it to graph partitions:
each worker returns $q$ local candidates, and a coordinator selects final seeds. \\
\midrule
Wang et al. 2023 & PaC-IM combines sketch compression and parallel seed
selection for large-scale influence maximization. It identifies space use and
sequential CELF-style seed selection as bottlenecks. &
Motivates why our course implementation focuses on runtime, spread-evaluation
count, and scalability. We do not need to reproduce PaC-IM, but can cite it as
evidence that parallel seed selection remains important. \\
\midrule
Chen et al. 2024 & Studies scalable distributed algorithms for size-constrained
monotone submodular maximization in MapReduce and adaptive-complexity models. &
Gives modern theoretical context for distributed submodular maximization. Useful
for framing our simplified GreeDI-style approach as a course-scale version of
current research. \\
\midrule
Barik et al. 2025 & GreediRIS combines reverse influence sampling with
distributed streaming maximum coverage and communication reduction techniques. &
Most relevant recent paper for communication constraints. It motivates reporting
candidate communication volume $mq$ and explaining the tradeoff between quality
and communication. \\
\bottomrule
\end{longtable}

\section{Core Definitions and Results to Reuse}

\subsection{Influence Maximization}

Given a graph $G=(V,E)$, edge propagation probabilities, a diffusion model $M$,
and a seed budget $k$, influence maximization asks for
\[
    \max_{S\subseteq V,\ |S|\le k} \sigma(S),
\]
where $\sigma(S)$ is the expected number of activated vertices when the seed set
$S$ is initially active.

For our project, the \textbf{Independent Cascade model} is the main diffusion
model.  When a node $u$ first becomes active, it gets one independent chance to
activate each currently inactive out-neighbor $v$ with probability $p_{uv}$.
The process stops when no further activation occurs.  A Monte Carlo estimator
can approximate $\sigma(S)$ by averaging the final active-set size over $R$
simulation runs.

\subsection{Submodularity}

A set function $f:2^V\rightarrow \mathbb{R}$ is submodular if for any
$A\subseteq B\subseteq V$ and $v\notin B$,
\[
    f(A\cup\{v\})-f(A)\ge f(B\cup\{v\})-f(B).
\]
This is the diminishing-returns property. For influence maximization under IC
and LT, $\sigma(S)$ is monotone and submodular. This directly justifies
greedy seed selection:
\[
    v^\star=\arg\max_{v\in V\setminus S}
    \left(\sigma(S\cup\{v\})-\sigma(S)\right).
\]

\subsection{Greedy Approximation Guarantee}

Nemhauser, Wolsey, and Fisher prove that for a nonnegative, monotone submodular
function under a cardinality constraint, the greedy algorithm returns a set $S$
with
\[
    f(S)\ge (1-1/e) f(S^\star),
\]
up to the finite-$k$ form $1-\left(\frac{k-1}{k}\right)^k$. Kempe et al. apply
this theorem to influence maximization after proving that the influence spread
function is monotone submodular. This is the main theory bridge in our report:
modern network diffusion becomes a classical submodular maximization problem.

\section{Paper-by-Paper Extraction}

\subsection{Nemhauser, Wolsey, and Fisher (1978)}

\textbf{Helpful content.}
\begin{itemize}
\item The paper studies maximizing a submodular set function with a cardinality
constraint, exactly the abstract optimization form behind influence maximization.
\item It treats submodularity as a discrete analogue of concavity: marginal
benefit decreases as the selected set grows.
\item For nondecreasing $z(S)$ with $z(\emptyset)=0$, greedy obtains a worst-case
ratio approaching $(e-1)/e=1-1/e$.
\end{itemize}

\textbf{How to use it.} In the theory section, state that the influence objective
is an instance of monotone submodular maximization under a cardinality
constraint. Then cite Nemhauser et al. for the generic greedy guarantee before
citing Kempe et al. for the influence-specific proof of submodularity.

\subsection{Kempe, Kleinberg, and Tardos (2003)}

\textbf{Helpful content.}
\begin{itemize}
\item The paper formulates the viral-marketing seed-selection problem as
influence maximization: choose $k$ initial adopters to maximize expected final
adoption.
\item It defines two classical diffusion models, IC and LT. Our project should
use IC as the primary implementation target because it is easy to simulate and
matches the proposal.
\item Influence maximization is NP-hard under IC and LT, so exact optimization
is not expected in a course implementation.
\item The expected influence spread $\sigma(S)$ is monotone and submodular under
IC and LT. Therefore greedy has a $(1-1/e-\epsilon)$-style approximation when
spread values are estimated sufficiently accurately.
\item Experiments compare greedy against degree and centrality heuristics and
show that heuristics can waste seeds on overlapping neighborhoods. This supports
our baselines: random, degree, and PageRank/centrality.
\end{itemize}

\textbf{Implementation takeaway.} Use Monte Carlo IC simulation for
$\widehat{\sigma}(S)$. The naive greedy cost is high because for each of $k$
rounds, many candidate nodes need marginal-gain estimates. This motivates CELF
and distributed candidate filtering.

\subsection{Leskovec et al. (2007)}

\textbf{Helpful content.}
\begin{itemize}
\item The paper studies outbreak detection but uses the same submodular
optimization structure: select a limited set of nodes to maximize a reward
function with diminishing returns.
\item It introduces CELF, a lazy greedy method using a priority queue of cached
marginal gains.
\item The key observation is that, by submodularity, a previously computed
marginal gain is an upper bound after the selected set grows. A node only needs
reevaluation when it reaches the top of the queue.
\item CELF preserves the greedy choice rule but can dramatically reduce the
number of expensive objective evaluations.
\end{itemize}

\textbf{Implementation takeaway.} Maintain entries $(v,\widehat{\Delta}_v,t_v)$,
where $t_v$ records the seed-set size when $\widehat{\Delta}_v$ was last
computed. Pop the largest cached gain. If it is current, select it; otherwise
recompute its marginal gain under the current seed set and push it back.
Report both runtime and the number of spread evaluations to show the benefit of
lazy evaluation.

\subsection{Mirzasoleiman et al. (2013)}

\textbf{Helpful content.}
\begin{itemize}
\item The paper proposes GreeDI, a simple distributed protocol for monotone
submodular maximization.
\item The protocol partitions the ground set across machines, runs greedy
locally, sends local selected elements to a coordinator, and runs greedy again
on the union of local candidates.
\item The main advantage is low communication: instead of sending the whole
dataset, each worker sends only a small candidate set.
\item The paper explains why distributed greedy is attractive but not identical
to centralized greedy: partitions may select redundant candidates or miss nodes
whose value is mainly global.
\end{itemize}

\textbf{Adaptation to our project.}
\[
    V=V_1\cup\cdots\cup V_m,\qquad
    C_i=\text{LocalGreedy}(G_i,q),\qquad
    C=\bigcup_{i=1}^{m}C_i,
\]
then the coordinator returns
\[
    S_{\mathrm{dist}}=\text{Greedy}(G,C,k).
\]
Our experiments should vary $m$ and $q$, then report
\[
    \rho=\frac{\sigma(S_{\mathrm{dist}})}{\sigma(S_{\mathrm{cent}})}
\]
and the communication proxy $mq$.

\subsection{Wang, Ding, Gu, and Sun (2023)}

\textbf{Helpful content.}
\begin{itemize}
\item The paper identifies two bottlenecks in large-scale influence
maximization: space-heavy sketches for marginal-gain estimation and sequential
seed selection.
\item PaC-IM combines sketch compression with parallel seed-selection data
structures. It is far beyond the scope of our reproduction, but it gives strong
motivation for measuring runtime and memory/scalability.
\item The paper distinguishes Monte Carlo simulation from sketch-based methods:
sketches reuse sampled graphs, often needing fewer samples than independent
simulation per evaluation.
\item It notes that CELF-style laziness helps reduce evaluations but can be
inherently sequential, which motivates parallel or distributed variants.
\end{itemize}

\textbf{How to use it.} Cite this paper in the related-work section as evidence
that scalable influence maximization is still an active systems problem. In our
scope section, explicitly say we implement Monte Carlo greedy/CELF and a simple
distributed greedy, while treating PaC-IM as advanced context.

\subsection{Chen, Dey, and Kuhnle (2024)}

\textbf{Helpful content.}
\begin{itemize}
\item The paper studies size-constrained monotone submodular maximization in
MapReduce and adaptive-complexity models.
\item It frames standard greedy as accurate but hard to scale because it needs
centralized access and many adaptive rounds.
\item It connects distributed algorithms with parallelizable low-adaptivity
algorithms, which is useful for explaining why distributed influence
maximization is not only about graph partitioning but also about reducing
sequential dependence.
\item It provides modern context for algorithms that keep approximation
guarantees while reducing communication or adaptivity.
\end{itemize}

\textbf{How to use it.} Use this paper to strengthen the background paragraph on
distributed submodular maximization. It is less necessary for implementation
details, but valuable for positioning our simplified GreeDI experiment as a
course-scale instance of a broader research direction.

\subsection{Barik et al. (2025)}

\textbf{Helpful content.}
\begin{itemize}
\item GreediRIS targets distributed influence maximization using reverse
influence sampling (RIS).
\item RIS converts seed selection into a maximum-$k$-coverage problem over
random reverse reachable sets. A seed is good if it covers many sampled reverse
reachable sets.
\item The paper focuses on the seed-selection bottleneck in distributed RIS
systems and uses a distributed streaming max-cover approach.
\item It explicitly studies communication bottlenecks and introduces truncation
as a knob that trades quality for lower communication.
\end{itemize}

\textbf{How to use it.} This is the strongest recent support for our
communication-constrained angle. We can cite it when defining $q$ as a local
candidate budget and when explaining why increasing $q$ should improve quality
but increase communication.

\section{Recommended Report Integration}

\subsection{Stage 1: Background and Related Work}

Use Kempe et al. for the application background and diffusion model. Use
Nemhauser et al. for submodular maximization. Use Leskovec et al. for lazy
evaluation. Use Mirzasoleiman et al., Chen et al., Wang et al., and Barik et al.
to motivate distributed and scalable variants.

\subsection{Stage 2: Representative Method}

The representative method should be \textbf{centralized greedy with IC Monte
Carlo estimation}, plus \textbf{CELF} as the main acceleration. The distributed
extension should be a GreeDI-style method rather than a full RIS system, because
it is easier to implement and directly matches the proposal.

\subsection{Stage 3--4: Implementation and Experiments}

Recommended algorithms:
\begin{itemize}
\item Random baseline.
\item Degree baseline.
\item PageRank or centrality baseline.
\item Monte Carlo greedy.
\item CELF/lazy greedy.
\item Distributed local greedy plus coordinator greedy.
\end{itemize}

Recommended metrics:
\begin{itemize}
\item estimated influence spread $\widehat{\sigma}(S)$;
\item runtime;
\item number of spread evaluations;
\item distributed quality ratio $\rho$;
\item communication proxy $mq$;
\item sensitivity to $m$, $q$, $k$, and propagation probability.
\end{itemize}

\section{Useful Text Snippets for the Final Report}

\textbf{Problem motivation.} Influence maximization turns viral marketing,
information diffusion, and intervention planning into a cardinality-constrained
combinatorial optimization problem on graphs. Its importance comes from the
fact that small seed sets can trigger large cascades, but selecting seeds
purely by local centrality may waste budget on overlapping neighborhoods.

\textbf{Theory bridge.} Under the IC model, the expected spread function is
monotone and submodular. Thus each new seed has diminishing marginal benefit as
the current seed set grows. This property connects influence maximization to
the classical Nemhauser-Wolsey-Fisher theorem, giving the greedy algorithm a
$(1-1/e)$ approximation guarantee.

\textbf{Why CELF.} The main cost of greedy influence maximization is repeated
spread estimation. CELF exploits submodularity by treating cached marginal
gains as upper bounds, recomputing a node only when it may still be the best
candidate. Therefore it preserves the greedy selection rule while reducing
objective evaluations.

\textbf{Why distributed greedy.} Centralized greedy assumes full access to the
graph and candidate set. In distributed settings, each worker can select local
candidates and send only a small set to a coordinator. This reduces
communication but may lose quality when important cross-partition influence is
not represented in the candidate pool.

\begin{thebibliography}{9}
\bibitem{nemhauser1978}
G. L. Nemhauser, L. A. Wolsey, and M. L. Fisher.
An analysis of approximations for maximizing submodular set functions - I.
\emph{Mathematical Programming}, 14:265--294, 1978.

\bibitem{kempe2003}
D. Kempe, J. Kleinberg, and E. Tardos.
Maximizing the spread of influence through a social network.
\emph{KDD}, 2003.

\bibitem{leskovec2007}
J. Leskovec, A. Krause, C. Guestrin, C. Faloutsos, J. VanBriesen, and N. Glance.
Cost-effective outbreak detection in networks.
\emph{KDD}, 2007.

\bibitem{mirzasoleiman2013}
B. Mirzasoleiman, A. Karbasi, R. Sarkar, and A. Krause.
Distributed submodular maximization: identifying representative elements in
massive data.
\emph{NeurIPS}, 2013.

\bibitem{wang2023}
L. Wang, X. Ding, Y. Gu, and Y. Sun.
Fast and space-efficient parallel algorithms for influence maximization.
\emph{PVLDB}, 17(3):400--413, 2023.

\bibitem{chen2024}
Y. Chen, T. Dey, and A. Kuhnle.
Scalable distributed algorithms for size-constrained submodular maximization in
the MapReduce and adaptive complexity models.
\emph{Journal of Artificial Intelligence Research}, 80:1575--1622, 2024.

\bibitem{barik2025}
R. Barik, W. Cappa, S. M. Ferdous, M. Minutoli, M. Halappanavar, and
A. Kalyanaraman.
GreediRIS: Scalable influence maximization using distributed streaming maximum
cover.
\emph{Journal of Parallel and Distributed Computing}, 198:105037, 2025.
\end{thebibliography}

\end{document}
